\reportheading{5.\quad 利用已有 FDS 结果安排检查顺序}
第 4 章以细网格 FDS 的最低净高判定方案可行性。本章用已有 FDS 数据检验两种排序策略：原布局按设计排序，改进布局在给定风量下按工况排序。两项回放分别进行。

\subhead{5.1\quad 设计排序：BO 与 GA}
在第 4.1 节的 16 个设计上，每次查询揭示一组已归档的 FDS 结果，记录首次找到最大净高余量的查询次序。固定顺序在第 16 次查询找到该方案，BO 在第 7 次找到，GA 在 20 个随机种子下于第 6 至 11 次找到。

\subhead{5.2\quad 工况排序：PINN 与回归代理模型}
在细网格回放中，固定顺序先查功率最低的风量，各风量内先查 OP-6、后查 OP-1，每种门态均按烟炱产率从低到高检查。这一顺序用 7 次检查确认最低功率可行方案。第 1 次在 OP-6、产率 0.02 下查到 0.96~m$^3$/s 越限，随后 6 次覆盖 1.12~m$^3$/s 的全部工况。保持相同的风量顺序，按预测净高排列工况时，PINN、直接回归、两者等权平均所需的检查次数分别为 7--11、7 和 7--9。两项回放的计数口径见图 6。

% Source: 05_Calculations/search_summary.csv (20 GA seeds) and
% 09_PINN_Layer/power_replay_summary.csv (refined mesh, scenario-only order).
\reportfigure{\begin{tikzpicture}
\pgfplotsset{replayaxis/.style={width=.48\linewidth,height=3.75cm,xmin=0,xmax=18,
 xtick={0,4,8,12,16},xmajorgrids,axis lines=left,
 tick label style={font=\sffamily\footnotesize},title style={font=\sffamily\small\bfseries,color=ink}}}
\begin{axis}[replayaxis,name=designrank,ymin=.5,ymax=3.5,ytick={1,2,3},
 yticklabels={GA,BO,固定},title={原布局：查询设计},xlabel={首次找到最大余量的查询数}]
\addplot[teal,thick,mark=*,mark size=2pt] coordinates {(6,1)(11,1)};
\addplot[teal,only marks,mark=*,mark size=2pt] coordinates {(7,2)};
\addplot[ink,only marks,mark=square*,mark size=2pt] coordinates {(16,3)};
\end{axis}
\begin{axis}[replayaxis,at={(designrank.east)},anchor=west,xshift=1.52cm,
 ymin=.5,ymax=4.5,ytick={1,2,3,4},yticklabels={混合,回归,PINN,固定},title={改进布局：查询工况},xlabel={确认最低功率方案的检查数}]
\addplot[rust,thick,mark=*,mark size=2pt] coordinates {(7,1)(9,1)};
\addplot[rust,only marks,mark=*,mark size=2pt] coordinates {(7,2)};
\addplot[rust,thick,mark=*,mark size=2pt] coordinates {(7,3)(11,3)};
\addplot[ink,only marks,mark=square*,mark size=2pt] coordinates {(7,4)};
\end{axis}
\end{tikzpicture}}{图 6.\quad 两项独立的 FDS 回放。左图的线段为 GA 在 20 个随机种子下的查询次数范围；右图的线段为缺少归档结果的工况取不同结局时的检查次数范围。}


\subhead{5.3\quad 最低净高的预测与 FDS 对照}
预筛数据包含 OP-1 和 OP-6 的 27 次完整 FDS 计算，覆盖 26 个物理设计。每次测试留出一个完整设计及其所有网格变体。修正后的 PINN、直接回归及二者等权平均对最低烟层高度的平均绝对误差分别为 6.97、7.39 和 6.70~cm。以 FDS 为参照，三种模型在 8 个越限算例中的漏判数依次为 7、4、4，在 19 个通过算例中的误报数依次为 0、1、0。按整类补风口布局留出时，各模型的误差和分类计数见表 F2。

图 5 的三个空心点是运行前冻结的 PINN 预测。对 OP-1、烟炱产率 0.05 和 0.10，PINN 分别低估最低净高 5.1 和 34.8~mm，混合模型相应低估 1.7 和 5.2~cm。对控制设计的 OP-6、产率 0.05 工况，PINN 单独预测为 0.5766~m，混合预测为 0.5576~m，FDS 则为 0.4358~m。相对于 FDS，PINN 与混合模型的预测分别高出 140.8 和 121.8~mm。该工况的 FDS 净高余量为 55.8~mm，六组复核结果见表 D3。

代理模型对整条曲线和最低净高的误差衡量不同。在采用物理信息极限学习机~\cite{pielm} 的整组留出时序试验中，PINN 在 16 个原布局和 6 个抬高补风口算例上的曲线平均绝对误差分别为 4.55 和 1.68~cm。附录 E 对照八个测点的时序，并按越限时段评价短时事件。
